A student used a telescope of diameter $D=100\,\mathrm{cm}$ to observe a distant star whose absolute magnitude is $-0.5$. If we consider the limiting magnitude of the eye as 6 and the diameter of the pupil $d$ to be $7\,\mathrm{mm}$, then answer the following questions.
\begin{enumerate}
\item If no interstellar medium (ISM) is present, find the maximum distance to which the student can detect this star.
\item In reality, an ISM of extinction factor $0.05\,\mathrm{mag\,kpc^{-1}}$ is present between the star and the Earth. Find the maximum distance to which the student can detect the star.
\end{enumerate}
\textbf{Note.} When an equation cannot be solved analytically, we use iterative methods to find an approximate numerical solution. As an example, suppose you want to find out the value of $x=x_{\mathrm{fin}}$, for which $x_{\mathrm{fin}}=f(x_{\mathrm{fin}})$. We start with some value $x=x_0$ and proceed as follows.
\begin{itemize}
\item Step 1. Find $f(x_0)$ and call it $x_1$.
\item Step 2. Find $f(x_1)$ and call it $x_2$.
\item Step 3. Find $f(x_2)$ and call it $x_3$.
\end{itemize}
Keep repeating this iteration process till $x_n$ is almost the same as $x_{n-1}$. That is your desired $x_{\mathrm{fin}}$ value.
